% A.57 full independent root and DE audits PASS; see notes/fgh_A57_audit.md.
\subsection{The rank-one case and all sufficiently large coefficient primes}

\begin{statement}{A.57}{Proved}{Potential good curve reduction excludes every noncuspidal pro-p normalizer}
Let \(Y\) be an affine hyperbolic curve over a number field, with
smooth proper compactification \(C\) of genus at least two.
If \(C\) has potentially good curve reduction at a place above
the coefficient prime \(p\), every nontrivial closed procyclic
subgroup of \(\pi_1(Y_{\overline K})^{(p)}\) whose arithmetic
normalizer has open image is contained in cusp inertia.
There is no condition on the special \(p\)-rank, and \(p=2\)
is included.

Consequently, for every geometrically connected degree-\(d\)
finite etale cover \(Y\to E\setminus\{O\}\), where
\(E:y^2=x^3-x\), the normalizer conclusion holds at every
coefficient prime
\[
                              p>d.                     \tag{A.57.1}
\]
More generally it holds whenever its proper compactification
has potentially good curve reduction at a place above \(p\).
These are assertions about actual normalized procyclic subgroups;
no fixed-subgroup calculation or fixed-element lifting is asserted.
\end{statement}

\begin{proof}
The special-rank-zero case is A.46, and special rank at least
two is A.56. We prove the remaining rank-one case, including
the precise reduction fact required for arbitrary cyclic covers.

\paragraph{An unramified p-cover above a rank-one special curve.}
Let \(R\) be a complete discrete valuation ring of mixed
characteristic \((0,p)\), let \(\mathcal C/R\) be a smooth
proper curve with geometrically connected fibres, and suppose
\(f(\mathcal C_{\overline k})=1\). Let
\(q:C'_{K_R}\to C_{K_R}\) be a geometrically connected finite
etale Galois cover of \(p\)-power degree; the cyclic case will
suffice below. Raynaud's Theorem 1 gives potentially good
reduction of \(J_{C'}\) and a semistable source with a tree
as dual graph. Every component contracted to a point of
\(\mathcal C_k\) has crystalline slopes strictly between
zero and one \cite[Theorem 1, printed p.179]{fghRaynaud1990Pgroups}.
We claim that the good special Jacobian of \(C'\) has
\(p\)-rank exactly one.

First make a finite extension on which the source is semistable,
and then work over the completion of the maximal unramified
extension. Its residue field is algebraically closed. This
last extension only computes geometric components and ranks;
all good-reduction assertions already descend to finite
extensions. Re-normalize the smooth base in the extended
generic covering field, obtaining a finite normal map
\[
                   \mathcal T\longrightarrow\mathcal C.
                                                        \tag{A.57.2}
\]
Finiteness follows from excellence of a complete discrete
valuation ring. Its generic source remains geometrically
connected. Properness gives a finite \(R\)-algebra
\(H^0(\mathcal T,\mathcal O_{\mathcal T})\); restriction to
the dense generic fibre embeds this algebra in \(K_R\),
because the latter fibre is a smooth geometrically connected
proper curve. Every element is integral over \(R\), so normality
of \(R\) gives
\[
                  H^0(\mathcal T,\mathcal O_{\mathcal T})=R.
                                                        \tag{A.57.3}
\]
Stein factorization therefore makes the special fibre connected.
This argument applies after every finite extension and after
the indicated geometric residue extension: it proves the
geometric connectedness needed here, rather than only
connectedness over the original finite residue field.

If (A.57.2) is etale, its special fibre is already a connected
etale Galois \(p\)-cover of \(\mathcal C_k\). Otherwise
purity of the branch locus on the regular surface
\(\mathcal C\) says that its reduced branch support is
the entire special fibre. Indeed the generic map is etale,
and the only possible vertical prime divisor is the smooth
irreducible fibre \(\mathcal C_k\); there are no isolated
codimension-two branch points. This branch support is
geometrically unibranch.

Apply Raynaud's Proposition 1 to (A.57.2). Its part (i)
makes the reduced inverse image of this branch support
geometrically unibranch. In particular it is locally
irreducible. A locally irreducible noetherian scheme has
pairwise disjoint irreducible components, which are open
and closed. Connectedness from (A.57.3) therefore gives a
\emph{single} irreducible special component \(T_0\).
Its decomposition group is the entire Galois group \(G\).
Part (ii) makes its generic inertia \(I\) constant at all
points of \(T_0\); it is normal in \(G\). Put
\(\mathcal T'=\mathcal T/I\), and let \(T'_0\) be the
image of \(T_0\). The same proposition states that
\(\mathcal T'\to\mathcal C\) is etale on a neighborhood
of \(T'_0\), and that \(T'_0\to\mathcal C_k\) is an
etale Galois cover of group \(G/I\)
\cite[Proposition 1(i)--(ii), printed pp.179--180]{fghRaynaud1990Pgroups}.

Here the neighborhood assertion is sufficient globally:
\(T'_0\) is the \emph{whole} special fibre as a topological
space, because \(\mathcal T\) has only one special component
and the finite quotient is surjective. The quotient map is
also etale on the entire generic fibre. Every point is in
one of these fibres, so
\[
                  \mathcal T'\longrightarrow\mathcal C
                  \quad\hbox{is finite etale everywhere.}
                                                        \tag{A.57.4}
\]
Thus \(T'_0\) is smooth, proper and connected. The map
\(T_0\to T'_0\) is the homeomorphism in Proposition 1;
on function fields its extension is purely inseparable.
Equivalently the inertia quotient has removed the whole
separable residue extension. The normalization \(Z_0\)
of \(T_0\) is consequently a finite radicial cover of
\(T'_0\). Over the perfect residue field a finite purely
inseparable map of smooth curves is a power of relative
Frobenius followed by an isomorphism: their function fields
have \(p\)-basis of cardinality one. Hence
\[
                            f(Z_0)=f(T'_0).             \tag{A.57.5}
\]
Alternatively their etale sites agree under this universal
homeomorphism, giving the same \(H^1(-,\F_p)\).
No equality between the radicial degree and \(|I|\) is
needed; possible vertical multiplicities play no role.

The unramified Deuring--Shafarevich formula gives
\[
                  f(T'_0)-1
                    =|G/I|\bigl(f(\mathcal C_k)-1\bigr)=0.
                                                        \tag{A.57.6}
\]
For the precise unramified formula see
\cite[Corollary 1.8 and equation (1.8.2), printed pp.35--36]
{fghCrew1984}. In the etale case of (A.57.2), apply this
formula directly with \(G\), obtaining the same conclusion.

The semistable source has just this one noncontracted
component: its function field is the residue function field
of the unique vertical valuation of \(\mathcal T\).
Every other component is contracted to a point of the
smooth base. The tree assertion and strict slopes in
Raynaud's Theorem 1 therefore give
\[
           f(\overline J_{C'})=f(Z_0)+
                  \sum_{Z\ {
m contracted}}f(J_Z)=1.  \tag{A.57.7}
\]
Here the special Jacobian in a compact-type model is the
product of the component Jacobians, and strict slopes
imply zero \(p\)-rank. This proves the claim.

\paragraph{The new Prym has no rational arithmetic line.}
For any such cyclic cover of degree \(m=p^n\), put
\(P_q=\ker(q_*:J_{C'}\to J_C)^0\).
The identity \(q_*q^*=[m]\) gives the actual isogeny
\[
                           J_{C'}\sim J_C\times P_q.   \tag{A.57.8}
\]
All its factors have good reduction after a finite extension,
by Raynaud and the good reduction of \(J_C\). Special
\(p\)-rank is additive in products and invariant under
isogeny. Equations (A.57.7)--(A.57.8) show
\(f(\overline P_q)=1-1=0\).
Statement A.46 then excludes every one-dimensional
\(\Q_p\)-subquotient of \(V_pP_q\), even after arbitrary
finite global or local extension. In particular no new
homology line can be preserved by an open arithmetic group.
This assertion holds for \emph{every} geometrically connected
cyclic \(p^n\)-cover of the good rank-one curve; we have not
imposed a specialization-order condition on its Picard point.

\paragraph{Choosing an arbitrary detecting cyclic cover.}
Now suppose the proposed noncuspidal normalized subgroup
has generator \(u\). By A.11 its proper homology class is
\(v=p^r v_0\ne0\), with \(v_0\) primitive in the
symplectic lattice \(T=H_1(C_{\overline K},\Z_p)\).
Since \(g\ge2\), its mod-\(p\) perpendicular space has
dimension \(2g-1\ge3\). Choose
\[
       \overline w\in\overline v_0^\perp
                           \setminus\F_p\overline v_0.
                                                        \tag{A.57.9}
\]
Perfect integral duality lifts \(\overline w\) to a
primitive \(w\in v_0^\perp\), exactly as in A.56.
Here no condition on the special image of \(w\) is needed.
For \(n>r\), the primitive character
\(\eta_n:T\to\Z/p^n\), given by pairing with
\(w\bmod p^n\) after choosing a root-of-unity generator,
defines a geometrically connected degree-\(p^n\) etale
proper cyclic cover. It descends with its deck action
after a finite number-field extension. Its Picard point
and an appropriate line-bundle trivialization also give
this cover directly. By construction
\[
       \eta_n(v)=0,
       \qquad v\bmod p^n\notin
                    (\Z/p^n)\operatorname{PD}(\eta_n).
                                                        \tag{A.57.10}
\]
The latter assertion follows by dividing a hypothetical
equality \(p^rv_0=a w\pmod {p^n}\) by \(p^r\)
and reducing modulo \(p\); primitivity of \(w\)
first forces \(p^r\mid a\).

The same generator \(u\) lies in the actual covering
subgroup of the pro-\(p\) fundamental group because
\(\eta_n(v)=0\). This covering subgroup is the pro-\(p\)
group of the actual punctured cover \(q^{-1}Y\), since
the cover has \(p\)-power Galois degree. Intersect the
arithmetic normalizer with its arithmetic covering group,
after the finite descent extension. The intersection has
finite index in the original normalizer, hence still has
open arithmetic image. The proper integral class
\(\widetilde v\) of \(u\) on \(C'\) pushes to \(v\)
and spans a stable rational line. Its projection to the
Prym in (A.57.8) must vanish by the preceding paragraph.
Thus \(\widetilde v\) lies in the rational old homology,
which equals the deck-invariant rational subspace.
The integral homology is torsion-free, so it is also
integrally invariant under the deck group.

The actual integral Schreier calculation of A.55--A.56,
in a basis adapted to \(\eta_n(b_1)=1\), gives
\[
 \begin{split}
 q_*H_1(C',\Z_p)^{\operatorname{Deck}(q)}
   ={}&\Z_p a_1\oplus p^n\Z_p b_1\\
     &{}\oplus\bigoplus_{j=2}^g
                  (p^n\Z_p a_j\oplus p^n\Z_p b_j).
 \end{split}                                            \tag{A.57.11}
\]
Modulo \(p^n\) this is precisely the
\(\operatorname{PD}(\eta_n)\)-line. Its membership for
\(v=q_*\widetilde v\) contradicts (A.57.10).
The choice \(n>r\) is why this argument handles a
generator with divisible proper homology; no extraction
of a nonexistent root of \(u\) is used. This completes
the rank-one case and the first assertion.

\paragraph{Arbitrary covers of the punctured elliptic curve.}
A nonconstant proper map to \(E\) has source genus at
least one. Genus one is already A.40 at every coefficient
prime. For genus at least two the first assertion applies
whenever the stated potential good curve reduction holds.
If the cover has degree \(d\ge2\) and \(p>d\), then
\(p\) is odd and \((E,O)\) has good marked reduction.
Its geometric Galois closure embeds in \(S_d\), so has
order prime to \(p\). The prime-to-\(p\) specialization
and tame local model proved explicitly in A.56 give
potentially good reduction of the proper covering curve.
That argument imposed no congruence on \(p\); the
congruence in A.56 was only used for its earlier rank-one
case. Thus the present criterion proves (A.57.1) for
every \(p>d\). For \(d=1\), including \(p=2\), use
the original elliptic criterion or A.40 directly.
\end{proof}
